\section{Discussion and Limitation}
In this work, we revisit how spatial grounding should be injected into VLAs and identify a simple but overlooked finding: given an adequate 2D grounding signal, lifting it into 3D and injecting it directly into the action head is what unlocks both spatial and instance-level task generalization. Built on this insight, our lightweight, model-agnostic method delivers substantial gains on both GR00T-N1.6 and $\pi_{0.5}$ on LIBERO-PRO benchmark, and remains effective in real-world settings.

\textbf{Limitations.} Our method rests on three assumptions that limit its scope. First, we assume each sub-goal is associated with a target object or region, leaving object-agnostic motions outside the formulation; lifting 2D trajectories into 3D and injecting them through a similar mechanism is a possible extension. Second, we assume an adequate 2D grounding signal that identifies the target, so the task generalization we report is instance-level within pick-and-place: faithfully acting on this target under unseen instruction–scene pairings, not generalizing to new skills. And we use relatively simple grounding sources to isolate the effect of representation and injection; robust grounding for more complex scenes or instructions remains an open problem orthogonal to our contribution. Third, lifting 2D grounding into 3D requires an RGB-D camera with calibrated intrinsics and extrinsics; relaxing this, e.g., via monocular depth estimation, is left to future work.